\begin{examplebox}
  \textbf{\textcolor{CExample}{例1｜已知方向，直接求新式（基础）}}
  \textbf{题目：}将$y=2x+1$的图像向上平移$3$个单位，求新解析式。

  \textbf{最快三步：}
  \begin{enumerate}[leftmargin=*,itemsep=0.2em,topsep=0.2em]
    \item 看方向：\textbf{向上}，所以在原来的整个式子\textbf{外面加$3$}。
    \item 写新式：$y=(2x+1)+3=\boxed{2x+4}$。
    \item 验证：原来$x=0$时$y=1$；现在$x=0$时$y=4$，确实升高$3$。
  \end{enumerate}
  \begin{center}
  \begin{tikzpicture}[x=1.25cm,y=0.53cm,>=stealth]
    \draw[->,gray] (-0.85,0)--(2.15,0) node[right] {$x$};
    \draw[->,gray] (0,-0.45)--(0,6.55) node[above] {$y$};
    \draw[gray!75,thick] (-0.60,-0.20)--(1.48,3.96);
    \draw[CExample,very thick] (-0.60,2.80)--(1.17,6.34);
    \fill[gray!80!black] (0,1) circle (1.9pt) node[right] {$1$};
    \fill[CExample] (0,4) circle (1.9pt) node[right] {$4$};
    \draw[->,CExample,thick] (0.36,1.18)--(0.36,3.80) node[midway,right,xshift=55pt] {上移$3$};
    \node[right] at (1.8,4.2) {$y=2x+1$};
    \node[right,CExample] at (1.8,6.1) {$y=2x+4$};
  \end{tikzpicture}

  \end{center}
  \textbf{快用：}直接在原式后面加$3$即可；不必逐点重画。

\end{examplebox}

\begin{examplebox}
  \textbf{\textcolor{CExample}{例2｜已知新旧解析式，反推平移（逆向）}}
  \textbf{题目：}$y=x^2$变成$y=x^2-2$，图像怎样移动？

  \textbf{最快三步：}
  \begin{enumerate}[leftmargin=*,itemsep=0.2em,topsep=0.2em]
    \item 对照：新式$=$原式$-2$，\textbf{减的是整个函数值}。
    \item 判断：\textbf{向下平移$2$个单位}。
    \item 验证：最低点从$(0,0)$变到$(0,-2)$，横坐标没变。
  \end{enumerate}
  \begin{center}
  \begin{tikzpicture}[x=0.97cm,y=0.67cm,>=stealth]
    \draw[->,gray] (-2.15,0)--(2.30,0) node[right] {$x$};
    \draw[->,gray] (0,-2.55)--(0,3.15) node[above] {$y$};
    \draw[gray!75,thick,domain=-1.70:1.70,samples=55] plot (\x,{\x*\x});
    \draw[CExample,very thick,domain=-1.70:1.70,samples=55] plot (\x,{\x*\x-2});
    \fill (0,0) circle (1.8pt);
    \fill[CExample] (0,-2) circle (1.8pt);
    \draw[->,CExample,thick] (0.32,-0.22)--(0.32,-1.80);
    \node[anchor=west,CExample] at (2,-3.1) {下移$2$};
    \node[right] at (2.0,2.55) {$y=x^2$};
    \node[right,CExample] at (2.0,-1.2) {$y=x^2-2$};
    \node[left] at (0,-2) {$-2$};
  \end{tikzpicture}

  \end{center}
  \textbf{快用：}把两个解析式相减，若差是一个常数，就看正负；\textbf{正数上移，负数下移}。

\end{examplebox}

\begin{examplebox}
  \textbf{\textcolor{CExample}{例3｜连续平移，一次算完（综合）}}
  \textbf{题目：}$y=-x^2$先向上移$1$个单位，再向下移$3$个单位，求最终解析式和顶点。

  \textbf{最快三步：}
  \begin{enumerate}[leftmargin=*,itemsep=0.2em,topsep=0.2em]
    \item 先算总变化：$+1-3=-2$，说明\textbf{最后比原来低$2$}。
    \item 直接写：$y=-x^2+1-3=\boxed{-x^2-2}$。
    \item 原顶点$(0,0)$，先到$(0,1)$，最后到$\boxed{(0,-2)}$。
  \end{enumerate}
  \begin{center}
  \begin{tikzpicture}[x=1.02cm,y=0.65cm,>=stealth]
    \draw[->,gray] (-2.0,0)--(2.10,0) node[right] {$x$};
    \draw[->,gray] (0,-4.25)--(0,2.05) node[above] {$y$};
    \draw[gray!80,thick,domain=-1.52:1.52,samples=60] plot (\x,{-\x*\x});
    \draw[blue!70!black,thick,dashed,domain=-1.52:1.52,samples=60] plot (\x,{1-\x*\x});
    \draw[CExample,very thick,domain=-1.52:1.52,samples=60] plot (\x,{-2-\x*\x});
    \fill[blue!70!black] (0,1) circle (1.8pt);
    \fill[gray!80!black] (0,0) circle (1.8pt);
    \fill[CExample] (0,-2) circle (1.8pt);
    \draw[->,blue!70!black] (0.34,0.11)--(0.34,0.91) node[midway,right] {$+1$};
    \draw[->,CExample] (-0.34,0.89)--(-0.34,-1.88) node[midway,left] {$-3$};
    \node[right,blue!70!black] at (2.2,1.1) {第一次};
    \node[right,gray!85!black] at (2.2,-1.1) {原图};
    \node[right,CExample] at (2.2,-3.1) {最终};
  \end{tikzpicture}

  \end{center}
  \textbf{快用：}只涉及上下平移时，先把所有“向上”记作正数、“向下”记作负数，\textbf{合起来只算一次}。
\end{examplebox}
